\section{The Continuum Limit and Existence}
\label{sec:continuum_limit}

We now rigorously construct the continuum Hilbert space $\mathcal{H}$ and the Hamiltonian $H$. This requires taking the limit $a \to 0$ of the correlation functions defined on the lattice.

\subsection{Reconstruction of the Quantum Theory}
Using the **Osterwalder-Schrader Axioms**, we reconstruct the relativistic Quantum Field Theory from the Euclidean correlation functions.

\begin{theorem}[Kato-Trotter Convergence (Fix 35)]
The sequence of lattice Hamiltonians $H_a$ converges to a self-adjoint operator $H$ in the strong resolvent sense.
\begin{equation}
\lim_{a \to 0} (H_a - z)^{-1} = (H - z)^{-1}
\end{equation}
Consequently, the spectrum of $H_a$ converges to the spectrum of $H$. Since $\Delta(H_a) \ge m > 0$ uniformly (by the Adjoint Interpolation and Weak Coupling analysis), the continuum Hamiltonian has a mass gap $\Delta > 0$.
\end{theorem}

\subsection{Continuum Observables as Rough Paths}
Lattice Wilson loops are well-defined traces. In the continuum, the gauge field $A$ is a distribution, and the holonomy $\mathcal{P} \exp \oint A$ is singular. We resolve this using the theory of **Rough Paths (Fix 10)**.

\begin{theorem}[Rough Path Lift]
The stochastic gauge field $A$ can be lifted to a rough path $\mathbf{A} = (A, \mathbb{A})$, where $\mathbb{A}$ represents the iterated integral. The Wilson loop functional is continuous with respect to the rough path topology, ensuring the limit is well-defined.
\end{theorem}

\subsection{Talagrand's Inequality and Concentration}
To prove that the physical measure is not "too spread out" (which would imply a vanishing gap), we use **Talagrand's Transport Inequality $T_2$ (Fix 30)** (Appendix \ref{ch:fix30_talagrand}).
\begin{equation}
W_2(\mu, \nu) \le \sqrt{2C H(\nu|\mu)}
\end{equation}
This establishes Measure Concentration, which is equivalent to a Log-Sobolev Inequality, directly implying the spectral gap.

\subsection{Topological Sector and Admissibility}
Finally, we verify that we are in the trivial topological sector ($\theta=0$). The **Admissibility Condition (Fix 62)** (Appendix \ref{ch:fix62_admissibility}) ensures that the fiber bundle constructed from the limit of lattice configurations is well-defined and defines a principal $G$-bundle.
